Indeed, let @@M0@@ lie in @@M1@@. After replacing @@M2@@ by the algebraic closure of a purely transcendental extension, which is allowed by 1.6 (i), we may suppose that there are @@M3@@ elements in @@M4@@ which are algebraically independent over all the @@M5@@. Then the @@M6@@ successive ``generic'' hyperplane sections of @@M7@@, defined by means of these elements, are still special positive @@M8@@-cycles of degree @@M9@@.

Arguing by induction on @@M10@@, it therefore remains to note that, when @@M11@@, we have @@M12@@, and, when @@M13@@, we have @@M14@@. For this last point, we may suppose that @@M15@@ itself comes from an integral scheme @@M16@@, by the injection @@M17@@; but then @@M18@@, which is @@M19@@ since @@M20@@ is a field.

\sgaStableTarget{sga6.e13.cor.6-11}\textbf{Corollary 6.11.} (i) \emph{Let @@M21@@ be two integers, let @@M22@@ for @@M23@@ @@M24@@, and let @@M25@@ @@M26@@. Let @@M27@@ be an algebraically closed field, and @@M28@@ a special positive @@M29@@-cycle of dimension @@M30@@ and degree @@M31@@, with Hilbert polynomial @@M32@@. Then, for @@M33@@, we have}
@@M34@@
\emph{Moreover, @@M35@@ embeds as a subscheme of @@M36@@ with @@M37@@, defined by (at most) @@M38@@ equations of degree @@M39@@.}

(ii) (Chow) \emph{Let @@M40@@ be a noetherian scheme. As @@M41@@ ranges over the algebraically closed extensions of the @@M42@@, @@M43@@, the special positive @@M44@@-cycles of bounded dimension and degree form a ``bounded family'' (either with the @@M45@@ regarded as subschemes of a fixed @@M46@@, or in an evident abstract sense).}

